\scanpage{41}
\mspage{40}
% Reused typescript, upside down relative to the foliation.
\begin{center}-3.13-\end{center}
The key property of stable geometric points turns out to be the following:

\uline{Proposition 3.6}: Let $x$ be a geometric point of $(\mathbb{P}_n)^{m+1}$. Then $x$ is stable if andonly if there are integers $N\geq N_0>0$ and there are monomials $P_0,\ldots,P_m$ in the expressions $D_{\alpha_0,\ldots,\alpha_n}$ such that, for $0\leq i\leq m$:

(i) $P_i$ is a section of $L_0^N\otimes\cdots\otimes L_i^{N-N_0}\otimes\cdots\otimes L_m^N$

(ii) $P_i(x)\ne0$

\uline{Proof}: Let $E=(\mathbb{R})^{m+1}$ be Euclidean space of dimension $m+1$. To each $D_{\alpha_0,\ldots,\alpha_n}$ associate a point $d_{\alpha_0,\ldots,\alpha_n}\in E$ as follows:
\[
\begin{aligned}
d_{\alpha_0,\ldots,\alpha_n}&=(\epsilon_0,\epsilon_1,\ldots,\epsilon_m)\\
\epsilon_i&=0\text{ if }i\notin\{\alpha_0,\ldots,\alpha_n\}\\
\epsilon_i&=1\text{ if }i\in\{\alpha_0,\ldots,\alpha_n\}
\end{aligned}
\]
Then $d$ determines which invertible sheaf $D$ is a section of. More generally, to any monomial $P$ is the $D$'s, \unsure{associate the corresponding|associatethe corresponding} sum of the $d$'s. Now define $C_x$ to be the convex cone spanned by those points $d_{\alpha_0,\ldots,\alpha_n}$ such that $D_{\alpha_0,\ldots,\alpha_n}(x)\ne0$, i.e. such that the components $x^{(\alpha_0)},\ldots,x^{(\alpha_n)}$ do not lie in a hyperplane. Then it is easy to see that:
